
\documentclass{article}
\usepackage{booktabs}
\usepackage{graphicx}
\usepackage{amsmath}
\usepackage{microtype}
\usepackage[colorlinks=true, allcolors=black]{hyperref}
 
\title{Two-Stage Prompting Does Not Mitigate MCQ Positional Bias in LLMs}
\author{Karl Hanna\thanks{Code and data available at \url{https://github.com/cotenthusiast/two-prompt-research}}}
\date{March 2026}
 
\begin{document}
\maketitle
 
\begin{abstract}
Multiple-choice question benchmarks are widely used to evaluate large language models, but MCQ evaluation is known to be sensitive to the order and labelling of answer options. This paper proposes and evaluates a two-stage prompting approach as a model-agnostic alternative to existing debiasing methods. The model first generates a free-text answer without seeing the options, then maps that answer to the closest option in a second step. The hypothesis is that separating reasoning from option selection should reduce positional bias. This approach is tested across four conditions, baseline, two-stage, cyclic permutation, and two-stage with cyclic permutation, on 1,000 stratified MMLU questions using GPT-4.1 mini, Gemini 2.5 Flash, and Llama 3.1 8B. The results show that two-stage prompting does not reduce positional bias and consistently increases mean absolute deviation from the ground-truth answer-position distribution across all three models. End-to-end accuracy is lower or negligibly different from baseline in all cases, driven largely by unscorable outputs at the matching step. Cyclic permutation remains the stronger robustness baseline. These results also highlight the importance of strict end-to-end evaluation when assessing debiasing methods, as conditional accuracy alone can obscure the practical cost of unscorable outputs.
\end{abstract}
 
\section{Introduction}
Large language models are increasingly evaluated using multiple-choice question benchmarks such as MMLU. However, MCQ evaluation has known reliability issues. Studies have documented positional bias, where models favour certain option positions regardless of content, leading to evaluations that do not reliably reflect model knowledge \cite{pezeshkpour-hruschka-2024-large}.
 
Several approaches have been proposed to mitigate these issues, such as PriDe \cite{zheng2023large} and permutation. PriDe estimates positional bias from a small calibration set and adjusts prediction probabilities accordingly. However, it requires access to output logits, which closed-source models do not provide. Permutation avoids this restriction, but it is computationally expensive. Cyclic permutation requires $n$ API calls for $n$ options, while full permutation requires $n!$ calls.
 
We propose a two-stage prompting approach as an alternative. This approach separates the reasoning from option selection. We test this alongside cyclic permutation across three models on 1000 MMLU questions. Our hypothesis is that separating the reasoning step from option labels should reduce positional bias. An initial pilot on 100 questions suggested a possible benefit for Gemini 2.5 Flash, which motivated the larger benchmark.
 
Our results show that two-stage prompting does not reduce bias and in some cases even shifts positional bias rather than eliminating it. Cyclic permutation remains the most effective approach. We discuss what these results suggest about the use of two-stage prompting for MCQ evaluation.
 
\section{Related Work}
 
\paragraph{Multiple-choice brittleness and option-order sensitivity.}
A growing body of work argues that large language models are not robust multiple-choice reasoners, even when their benchmark accuracy looks strong. Zheng et al.~\cite{zheng2023large} show that LLMs can behave as unreliable multiple-choice selectors, which motivates methods that try to reduce sensitivity in the final selection step. Pezeshkpour and Hruschka~\cite{pezeshkpour-hruschka-2024-large} further show that model performance can change substantially when answer options are reordered, suggesting that multiple-choice evaluation may reflect sensitivity to superficial formatting choices rather than stable underlying knowledge. Related concerns are also raised by Wang et al.~\cite{wang2024llms}, who argue that models may succeed in MCQA by selecting the least incorrect option rather than identifying a clearly justified correct answer. More recent work continues this line by emphasizing that benchmark performance can also be distorted by hidden or disguised biases in option-based evaluation settings~\cite{nowak2026abcd}.
 
\paragraph{Debiasing and robustness-oriented prompting methods.}
Several recent approaches try to make multiple-choice evaluation more robust by modifying the prompting or selection procedure. PriDe~\cite{zheng2023large} is a central example, using perturbation-based debiasing to reduce sensitivity in option selection and to test whether a model's answer stays stable across changes in presentation. BiasPrompting~\cite{vu2025biasprompting} similarly treats bias as a central issue in MCQ answering and proposes a prompting-based intervention aimed at improving behavior under bias-sensitive conditions. Together, these methods treat robustness as a core requirement for reliable evaluation rather than a minor implementation detail. Our work is closest to this line, but differs in focusing on a two-stage generation-then-mapping procedure and testing whether separating answer generation from option selection improves robustness to MCQ-specific biases.
 
\paragraph{Open-style answering and answer matching.}
Another line of work questions whether forced-choice evaluation is the right interface for measuring model knowledge at all. Myrzakhan et al.~\cite{myrzakhan2024open} argue for moving from purely multiple-choice formats toward open-style questions in leaderboard-style evaluation, reflecting a broader concern that multiple-choice constraints can distort what is actually being measured. Chandak et al.~\cite{chandak2025answer} push this further by arguing that answer matching can outperform standard multiple-choice evaluation, suggesting that free-form responses may give a more faithful picture of model competence than letter selection alone. This perspective directly motivates two-stage methods: if the model is first allowed to answer in open form and only later mapped back to the provided options, the evaluation can test whether part of the observed MCQ brittleness comes from the selection interface rather than from the model's underlying ability to solve the problem.
 
\paragraph{Positioning of this work.}
This paper sits at the intersection of these strands. Prior work has shown that MCQ evaluation is vulnerable to option-order effects, selector bias, and other formatting-sensitive artifacts~\cite{zheng2023large,pezeshkpour-hruschka-2024-large,wang2024llms,nowak2026abcd}, while separate work has argued that open-form answering or answer matching may provide a better measurement interface~\cite{myrzakhan2024open,chandak2025answer}. Our contribution connects these ideas by testing whether a two-stage pipeline, first generating an answer and then mapping it to the provided options, improves robustness relative to direct multiple-choice selection and to perturbation-based debiasing baselines such as PriDe~\cite{zheng2023large}. In this sense, the paper is primarily about robustness under biased MCQ conditions, with faithfulness as a secondary motivation and raw accuracy as a supporting outcome rather than the main claim.
 
\section{Methodology}
MMLU is a benchmark dataset containing 14,000 questions across 57 subjects, covering a broad range of fields and difficulty levels. In this paper, 1,000 questions stratified across MMLU subjects were used.

Four conditions were evaluated across three LLM models. The \textbf{baseline} condition presents the question directly in multiple-choice format. \textbf{Cyclic permutation} asks the model the same question n times, cycling the answer options across all positions. After n responses are compiled (for n=4 options), each response is unpermuted to recover the original option and a majority vote is used to determine the final answer; ties are broken by defaulting to the original permutation.

\textbf{Two-stage prompting} first asks the model to generate a free-text response without the answer options visible, then presents the options alongside that response and asks the model to select the closest match. \textbf{Two-stage + cyclic permutation} combines both interventions: the free-text response is generated once, then the matching step is repeated four times with the options permuted across positions.

The three models evaluated are GPT-4.1 mini, Gemini 2.5 Flash, and Llama 3.1 8B (via Groq), representing a frontier closed-source model, an efficient closed-source model, and an open-weights model respectively, spanning different capability tiers and access patterns. A temperature of 0.0 was used across all models. A concurrency limit of 10 was applied to GPT-4.1 mini; Gemini 2.5 Flash and Llama 3.1 8B were limited to 2 concurrent requests to accommodate stricter rate limits.

The full prompt templates for each condition are provided in Appendix~\ref{appendix:prompts}. The key design decision in the free-text stage is that the question is presented without any answer options, ensuring that option labels cannot influence the generated response. The matching stage then receives both the free-text response and the original options and asks the model to select the closest match.

Two accuracy metrics are reported. End-to-end accuracy is defined as the number of correct answers divided by the total number of questions; unscorable outputs are counted as incorrect. Conditional accuracy is defined as the number of correct answers divided by the number of scorable outputs only. An output is considered unscorable if the matching step in two-stage prompting produces a response that cannot be parsed as A, B, C, or D, or if an API error occurs. To assess robustness, we compute the mean absolute deviation of the predicted answer-position distribution from the ground-truth answer-position distribution across all questions.
 
Additional configuration details, including random seed, maximum token limits etc., are defined in the project configuration files and are available in the repository. Gemini 2.5 Flash two-stage + cyclic results were truncated at n=196 due to provider spending cap exhaustion and are excluded from cross-condition comparisons.
 
\section{Results}

Table~\ref{tab:accuracy} shows end-to-end and conditional accuracy across all conditions. The main finding is straightforward: two-stage prompting does not improve over baseline and in most cases hurts performance. GPT-4.1 mini drops from 81.9\% to 78.6\%, and Gemini 2.5 Flash drops more sharply from 83.6\% to 65.5\%. Llama 3.1 8B is the only model where two-stage prompting does not hurt end-to-end accuracy, though the difference is negligible (+0.1pp). Cyclic permutation is the only method that consistently matches or beats baseline across all three models.

One important pattern in Table~\ref{tab:accuracy} is the gap between end-to-end and conditional accuracy. Conditional accuracy stays relatively stable across methods --- Gemini 2.5 Flash under two-stage prompting scores 87.3\% conditionally, not far from its 89.5\% baseline. The end-to-end drop to 65.5\% comes entirely from the 250 unscorable outputs at the matching step. This means the method is not making the model worse at answering questions --- it is just failing to produce a scorable answer for a large chunk of them, which still counts against it in practice.

\begin{table}[ht]
\centering
\caption{End-to-end and conditional accuracy (\%) by method and model. Gemini 2.5 Flash two-stage + cyclic results are truncated at $n=196$ due to provider spending cap exhaustion and should not be compared directly with other conditions.}
\label{tab:accuracy}
\begin{tabular}{lrrrrrr}
\toprule
  & \multicolumn{2}{c}{GPT-4.1-mini} & \multicolumn{2}{c}{Gemini-2.5-Flash} & \multicolumn{2}{c}{Llama-3.1-8B} \\
\cmidrule(lr){2-3}\cmidrule(lr){4-5}\cmidrule(lr){6-7}
Method & E2E & Cond. & E2E & Cond. & E2E & Cond. \\
\midrule
Baseline & 81.9 & 81.9 & 83.6 & 89.5 & 65.0 & 65.0 \\
Two-Stage & 78.6 & 80.0 & 65.5 & 87.3 & 65.1 & 65.3 \\
Cyclic Perm. & 81.6 & 81.6 & 86.3 & 88.9 & 65.4 & 65.4 \\
Two-Stage + Cyclic & 80.6 & 80.6 & --- & --- & 62.8 & 64.5 \\
\bottomrule
\end{tabular}
\end{table}

Table~\ref{tab:bias} shows positional bias measured as mean absolute deviation from the ground-truth answer-position distribution. Two-stage prompting increases MAD relative to baseline for every model: GPT-4.1 mini goes from 0.45 to 0.71, Gemini 2.5 Flash from 0.75 to 2.13, and Llama 3.1 8B from 2.55 to 3.14. So not only does two-stage prompting fail to reduce positional bias, it appears to make it worse. Cyclic permutation is the only method that brings MAD down, with the biggest improvement on Llama 3.1 8B where it drops from 2.55 to 0.65.

\begin{table}[ht]
\centering
\caption{Mean absolute deviation from ground-truth answer-position distribution (pp). Lower is better.}
\label{tab:bias}
\begin{tabular}{lrrr}
\toprule
Method & GPT-4.1-mini & Gemini-2.5-Flash & Llama-3.1-8B \\
\midrule
Baseline & 0.45 & 0.75 & 2.55 \\
Two-Stage & 0.71 & 2.13 & 3.14 \\
Cyclic Perm. & 0.50 & 0.99 & 0.65 \\
Two-Stage + Cyclic & 0.95 & --- & 2.93 \\
\bottomrule
\end{tabular}
\end{table}

Table~\ref{tab:failures} breaks down where unscorable outputs come from. For Gemini 2.5 Flash under two-stage prompting, 225 of the 250 unscorable outputs are parse failures at the matching step --- the model returned something that could not be read as A, B, C, or D. This is not a provider issue, it is the matching step itself failing. Llama 3.1 8B shows high API failure counts under cyclic permutation (349) and two-stage + cyclic (322), which were caused by Groq spending cap exhaustion during the run. These are recorded as failures rather than wrong answers, so scored outputs for those conditions remain at 1000 and 974 respectively. The Gemini two-stage + cyclic condition is excluded from analysis entirely, as 778 API failures left only 196 valid outputs.


\begin{table}[ht]
\centering
\caption{Provider failures, parse failures, and final unscorable questions per condition.}
\label{tab:failures}
\begin{tabular}{llrrrrr}
\toprule
Method & Model & Total & API Fail & Parse Fail & Unscorable & Scored \\
\midrule
Baseline & GPT-4.1-mini & 1000 & 0 & 0 & 0 & 1000 \\
Baseline & Gemini-2.5-Flash & 1000 & 8 & 58 & 66 & 934 \\
Baseline & Llama-3.1-8B & 1000 & 0 & 0 & 0 & 1000 \\
Two-Stage & GPT-4.1-mini & 1000 & 16 & 1 & 17 & 983 \\
Two-Stage & Gemini-2.5-Flash & 1000 & 25 & 225 & 250 & 750 \\
Two-Stage & Llama-3.1-8B & 1000 & 0 & 3 & 3 & 997 \\
Cyclic Perm. & GPT-4.1-mini & 1000 & 0 & 0 & 0 & 1000 \\
Cyclic Perm. & Gemini-2.5-Flash & 1000 & 6 & 28 & 29 & 971 \\
Cyclic Perm. & Llama-3.1-8B & 1000 & 349 & 0 & 0 & 1000 \\
Two-Stage + Cyclic & GPT-4.1-mini & 1000 & 2 & 0 & 0 & 1000 \\
Two-Stage + Cyclic & Gemini-2.5-Flash & 1000 & 778 & 28 & 804 & 196 \\
Two-Stage + Cyclic & Llama-3.1-8B & 1000 & 322 & 0 & 26 & 974 \\
\bottomrule
\end{tabular}
\end{table}
 
\section{Discussion}
Looking at these results, two-stage prompting clearly does not mitigate positional bias. It consistently increases bias across all three models, and end-to-end accuracy is either lower or negligibly different from baseline. The intuition behind the approach is reasonable --- in practice, people tend to read a question, form an answer on their own, and only then look at the choices to avoid anchoring on any particular option. This method tries to mirror that process, but it does not translate well to LLMs.

There are a few potential reasons for this. The bias may not disappear but simply move to the matching stage, where the model still has to pick from labelled options and remains susceptible to the same positional effects. The matching step also introduces a new failure surface that the baseline does not have --- outputs that seem reasonable but cannot be parsed as a final letter answer, which under strict evaluation count against the method. Some question types make this worse: a question like "which of the following countries is not in Europe?" cannot be answered without seeing the options first, so the free-text stage breaks down by design for these cases.

This study has a few limitations worth noting. Only three models were tested on a single benchmark, and only a subset of that benchmark was used. The evaluation is also limited to four-choice questions, so it is unclear whether the results generalise to settings with more or fewer options. Without access to model logits it is also hard to say with confidence what is happening internally during the matching step.

There are some directions worth exploring in future work. Better parsing, a baseline fallback for unscorable outputs, and allowing some temperature in the matching step could each address specific weaknesses found here. The matching prompt itself was also not varied --- different prompt designs for the matching stage may produce meaningfully different results. Whether these changes are enough to make two-stage prompting a viable robustness intervention is an open question.
 
\section{Conclusion}
This paper investigated whether two-stage prompting, which separates free-text answer generation from final option selection, can reduce positional bias in multiple-choice evaluation of large language models. Across three models and 1,000 MMLU questions, the answer is no. Two-stage prompting does not mitigate positional bias and in most cases increases it, with mean absolute deviation rising relative to baseline for every model tested.

The more important finding is where the damage shows up. Conditional accuracy stays relatively stable across conditions, which would suggest the method is largely harmless. The end-to-end numbers tell a different story: the matching step produces a substantial unscorable burden that counts against the method in practice. Evaluating only on scorable outputs would have obscured this entirely. Cyclic permutation remains the more reliable robustness intervention in this setting, despite its higher inference cost.

Two-stage prompting is not without potential. A small but consistent set of questions were answered correctly under two-stage that the baseline missed, suggesting there is some signal in the free-text generation step. A version of the method with improved parsing and a baseline fallback for unscorable outputs may be worth revisiting. Whether that is enough to make it a viable robustness intervention is left for future work.
 
\bibliographystyle{plain}
\bibliography{refs}
 
\appendix
 
\section{Prompt Templates}
\label{appendix:prompts}
 
The following prompts were used across all conditions. All models received identical prompts with no model-specific modifications.
 
\paragraph{Direct MCQ (Baseline and Cyclic).}
\begin{verbatim}
Answer the following multiple-choice question.
 
Question: {question}
 
Options:
A. {option_a}
B. {option_b}
C. {option_c}
D. {option_d}
 
Respond with only the letter.
\end{verbatim}
 
\paragraph{Free-Text Stage (Two-Stage and Two-Stage + Cyclic).}
\begin{verbatim}
Answer the following question based on your knowledge.
 
Question: {question}
\end{verbatim}
 
\paragraph{Option Matching Stage (Two-Stage and Two-Stage + Cyclic).}
\begin{verbatim}
You are given a question, a reference answer, and four options.
 
Question: {question}
 
Reference answer: {free_text}
 
Options:
A. {option_a}
B. {option_b}
C. {option_c}
D. {option_d}
 
Select the option that best matches the reference answer in the
context of the question.
If the reference answer is imperfect or incomplete, choose the
closest option.
Respond with only the letter.
\end{verbatim}
 
\end{document}